\section{Predictive usefulness and the justification of complexity}
\label{sec:utility-question}

The question ``when is a forecasting representation useful?'' depends on
what information is retained, which predictor is fitted, and what resources
the complete procedure consumes. We distinguish these objects before
comparing KAHM representations. Let $X$ be an available causal history,
$Y$ its fixed future target, and $\Gamma_{\mathcal D}$ a representation
fitted using development data $\mathcal D$. A forecasting procedure includes
the representation, decoder, regularization, and selection rule. Conditional
on its fitted state, it returns a fixed forecast $\widehat f$ for a new
evaluation unit. The studies below specify their own target, origin law,
aggregation, and evaluation unit; their risks are not interchangeable.

\paragraph{Information retention and finite-decoder access.}
A representation may discard distinctions needed to predict $Y$. This is
different from making those distinctions difficult for the chosen decoder
to use. In particular, appending a deterministic feature map to the raw
history gives $(X,\Gamma_{\mathcal D}(X))$, from which $X$ is exactly
recoverable. Conditional on the fitted state, this preserves the
raw-history Bayes information; it does not create additional information
about the future (Proposition~\ref{prop:nf-raw-preservation}). An improvement
from raw augmentation must concern the finite fitting procedure, its
regularization, or its numerical behavior. Conversely, an information-
preserving augmentation can worsen that procedure's generalization.

\paragraph{Predictive structure relative to a fixed baseline.}
There is a useful oracle distinction before selecting a finite decoder.
Condition on the fitted state, fix a raw forecast $f_0$, and let
$e=Y-f_0(X)$ and $Z=\Gamma_{\mathcal D}(X)$, with
$\E\|e\|_2^2<\infty$. Write $R(f)=\E\|Y-f(X)\|_2^2$ for squared-error
risk under this conditional evaluation law. The standard conditional-expectation projection
identity gives the best square-integrable correction
$q^*(Z)=\E[e\mid Z]$ and
\begin{equation}\label{eq:utility-oracle}
 R(f_0)-\inf_{q\in L_2(P_Z)}R(f_0+q\circ\Gamma_{\mathcal D})
 =\E\|\E[e\mid Z]\|_2^2.
\end{equation}
Thus the representation contains useful mean-predictive structure for
this baseline's residual precisely when that conditional mean is nonzero
with positive probability. This is a standard projection fact, not a new
KAHM theorem (Appendix~\ref{app:utility-projection}). It is relative to
the specified baseline and evaluation law: unpredictable residual
variation does not become useful structure merely because baseline MSE
is large. A deterministic $Z$ can make existing structure easier for a
restricted procedure to use without adding information to the raw history.
The oracle ranges over all square-integrable corrections; neither its
gain nor the conditional mean is estimated by our finite experiment.
Whether a particular learned correction accesses that structure is a
separate empirical question.

\paragraph{An operational utility comparison.}
For two frozen forecasts $\widehat f_0$ and $\widehat f_s$, evaluated on the
same unit $E$ and squared-loss aggregation $\ell$, define
\begin{equation}\label{eq:utility-delta}
 \Delta_s=\E\!\left[\ell(\widehat f_0;E)-\ell(\widehat f_s;E)
             \mid\mathcal D\right].
\end{equation}
Positive $\Delta_s$ means the richer fitted forecast improves expected
loss under that evaluation law. A practically useful improvement must also
be compared with a declared margin and measured resource requirements.
Selection is part of the procedure: a rule pays for the features, probes,
and candidate fits needed to make its choice even when it retains the
baseline. Evaluation of that rule uses outcomes unavailable to its fitting
and selection. Repetitions of training and independently sampled evaluation
units assess different forms of variability.

The procedure must also return a forecast on the declared task. A model
whose candidate class contains the baseline can still fail during feature
fitting or search before selecting that baseline. Without a prespecified
failure fallback, its full-roster risk is undefined. We therefore retain
fit coverage as part of the utility assessment, rather than treating
missing predictions as favorable error outcomes.

For an application that specifies a minimum useful improvement $\delta$
and a resource budget $B$, the operational question is whether the evidence
supports $\Delta_s\geq\delta$ at a complete cost $C_s\leq B$, compared with
the available alternatives. A large baseline error alone is insufficient:
it can include unpredictable variation as well as structure accessible to
a correction. The present study reports error--cost profiles and tests a
fixed validation rule. It does not assume a universal deployment budget or
identify a universal exchange rate between prediction error and computation.

\paragraph{What an empirical cost decision establishes.}
An observed accuracy--cost advantage concerns the fitted procedures being
compared. It does not show that every cheaper representation or decoder
fails an accuracy requirement. The stronger necessity statement requires
lower bounds that apply to the entire class, a declared sensitivity budget, and the
sampling assumptions of the related certificate framework
\citep{Anonymous2026Certificates}. Here we study empirical predictive
utility and measured cost. We neither estimate unrestricted Bayes risk on
natural data nor certify minimum feasible representation complexity.

\paragraph{Preserving a useful forecasting pathway.}
Raw-history retention alone need not retain the original finite decoder
class. In the original natural study, concatenating features changes
Gaussian distances and the Nystr\"om basis (Section~\ref{sec:natural}).
The extension therefore examines an additive forecast
\begin{equation}\label{eq:utility-additive}
 \widehat f_{g,\gamma}(x)
       =\widehat f_0(x)+\gamma\widehat q_g(x),\qquad 0\leq\gamma\leq1,
\end{equation}
where the raw forecast $\widehat f_0$ remains fixed while a correction
$\widehat q_g$ is fitted using KAHM, affine, or locality features. Setting
$\gamma=0$ recovers the baseline exactly. Residual fitting is established
\citep{Buhlmann2003L2}; the question is which geometric correction, if any,
is useful under the specified data and computation budgets.

For fixed $\widehat f_0$ and $\widehat q_g$, an elementary squared-loss
identity gives
\begin{equation}\label{eq:utility-alignment}
 R(\widehat f_{g,\gamma})-R(\widehat f_0)
 =\gamma^2\E\|\widehat q_g(X)\|_2^2
 -2\gamma\E\langle Y-\widehat f_0(X),\widehat q_g(X)\rangle.
\end{equation}
All expectations condition on the frozen fitted state. For a nonzero
correction, positive residual alignment is necessary and permits a
sufficiently small positive mixing weight to improve the population risk.
The identity is a motivation, not a new risk certificate: its population
quantities are unknown during fitting, a finite weight grid can miss the
improving interval, and validation can select a worse forecast.
Appendix~\ref{app:utility-alignment} provides the derivation and the
distinction between feasible baseline recovery and fitted generalization.
